% GENERATED FILE. Edit canonical lessons, then run npm run generate:books.
% canonical-source-hash: fnv1a64:68a4399243b925a1
% canonical-lessons: ZH-C126-changchang, ZH-C126-youshihou, ZH-C126-yiqian, ZH-W126-yi-use, ZH-C126-gangcai, ZH-C126-yijing

\chapter{Often and Sometimes}
\label{ch:zh-often-and-sometimes}
\hlchaptermodality{126}

\begin{chapteropening}
\textbf{By the end of this chapter:} \emph{I can say how often I do things, what I did before and what I did a moment ago.}

Complete the last lesson of chapter 126: I can say how often I do things, what I did before and what I did a moment ago.
\end{chapteropening}

\section[chángcháng]{chángcháng — often}

\label{lesson:ZH-C126-changchang}



\begin{quote}
Before the new one: say the Mandarin for anxious, then the Mandarin for upset.
\end{quote}

\subsection*{You'll want to know: chángcháng}
\emph{chángcháng} — \textquotedblleft{}often\textquotedblright{}.

\begin{quote}
\emph{wǒ chángcháng hē chá} — I often drink tea.
\end{quote}

Its characters come later; for now it is a word for the ear.

\begin{cousinweb}
\emph{chángcháng} says \emph{cháng} twice. It is inherited from Middle Chinese \emph{dzyang} (Guangyun, in Baxter's transcription), which glosses it as a length of two fathoms and as an official title. The sense 'usual' is not given there, so how it came to mean often is not settled.
\end{cousinweb}

\begin{grammarlens}[title={What you've built}]
One more word for times of day.
\end{grammarlens}

\subsection*{Your turn}
\begin{itemize}
\raggedright
  \item \emph{Say it:} \emph{chángcháng}
  \item \emph{Say it:} \emph{chángcháng}, once more
  \item \emph{Say it:} \emph{wǒ chángcháng hē chá} — I often drink tea.
  \item \emph{Recall:} say the Mandarin for anxious, then the Mandarin for upset, then say \emph{chángcháng} again
\end{itemize}

\begin{tcolorbox}[breakable,colback=teal!4,colframe=teal!35!black,title={Before you move on}]
What does chángcháng mean? (\textquotedblleft{}Often\textquotedblright{}.) Say it once more.
\end{tcolorbox}
\section[yǒushíhou]{yǒushíhou — sometimes}

\label{lesson:ZH-C126-youshihou}



\begin{quote}
Before the new one: say the Mandarin for upset, then the Mandarin for often.
\end{quote}

\subsection*{You'll want to know: yǒushíhou}
\emph{yǒushíhou} — \textquotedblleft{}sometimes\textquotedblright{}.

\begin{quote}
\emph{wǒ yǒushíhou qù Běijīng} — Sometimes I go to Beijing.
\end{quote}

Its characters come later; for now it is a word for the ear.

\begin{cousinweb}
\emph{yǒushíhou} is built from \zh{有} 'there is' and \emph{shíhou} 'time', all three parts inherited from Middle Chinese, \emph{hjuwX}, \emph{dzyi} and \emph{huwH} (Guangyun, in Baxter's transcription): there are times.
\end{cousinweb}

\begin{grammarlens}[title={What you've built}]
One more word for times of day.
\end{grammarlens}

\subsection*{Your turn}
\begin{itemize}
\raggedright
  \item \emph{Say it:} \emph{yǒushíhou}
  \item \emph{Say it:} \emph{yǒushíhou}, once more
  \item \emph{Say it:} \emph{wǒ yǒushíhou qù Běijīng} — Sometimes I go to Beijing.
  \item \emph{Recall:} say the Mandarin for upset, then the Mandarin for often, then say \emph{yǒushíhou} again
\end{itemize}

\begin{tcolorbox}[breakable,colback=teal!4,colframe=teal!35!black,title={Before you move on}]
What does yǒushíhou mean? (\textquotedblleft{}Sometimes\textquotedblright{}.) Say it once more.
\end{tcolorbox}
\section[\texorpdfstring{\zh{以前} (yǐqián)}{以前 (yǐqián)}]{\zh{以前} (yǐqián) — before; in the past}

\label{lesson:ZH-C126-yiqian}



\begin{quote}
Before the new one: say the Mandarin for often, then the Mandarin for sometimes.
\end{quote}

\subsection*{You'll want to know: \zh{以前}}
\textbf{\zh{以前}} — \emph{yǐqián} — \textquotedblleft{}before; in the past\textquotedblright{}.

\begin{quote}
\textbf{\zh{以前我不会说汉语}} \emph{yǐqián wǒ bú huì shuō hànyǔ} — I couldn't speak Chinese before.
\end{quote}

The new character, in pieces: \zh{以} — four strokes: a vertical that flicks up at its foot, a dot beside it, then a long stroke falling to the left and a dot at the lower right. The next lesson puts it in your hand.

\begin{cousinweb}
\textbf{\zh{以前}} (\emph{yǐqián}) is built from \zh{以}, glossed 'to use' and 'with', and \zh{前} 'before', both inherited from Middle Chinese, \emph{yiX} and \emph{dzen} (Guangyun, in Baxter's transcription).
\end{cousinweb}

\begin{grammarlens}[title={What you've built}]
One more word for times of day.
\end{grammarlens}

\subsection*{Your turn}
\begin{itemize}
\raggedright
  \item \emph{Say it:} \emph{yǐqián}
  \item \emph{Say it:} \emph{yǐqián}, once more
  \item \emph{Say it:} \emph{yǐqián wǒ bú huì shuō hànyǔ} — I couldn't speak Chinese before.
  \item \emph{Recall:} say the Mandarin for often, then the Mandarin for sometimes, then say \emph{yǐqián} again
\end{itemize}

\begin{tcolorbox}[breakable,colback=teal!4,colframe=teal!35!black,title={Before you move on}]
What does \zh{以前} mean? (\textquotedblleft{}Before; in the past\textquotedblright{}.) Say it once more.
\end{tcolorbox}
\section[\texorpdfstring{\zh{以} (yǐ)}{以 (yǐ)}]{\zh{以} — four strokes in two pairs}

\label{lesson:ZH-W126-yi-use}



\begin{quote}
The word first, then the character inside it.

\begin{itemize}
\raggedright
  \item \emph{Recall:} say \emph{yǐqián} — \textbf{R1}, one lesson back, before it decays
\end{itemize}

\begin{quote}
\zh{前} \zh{会}
\end{quote}

\emph{Read it:} each one aloud, with nothing said first — \emph{qián}, \emph{huì}
\end{quote}

\subsection*{Script — four strokes in two pairs}
\hlblockfigure{\detokenize{figures/ZH-W126-yi-use-filmstrip.pdf}}{How \zh{以} is written, stroke by stroke}

\begin{quote}
\zh{以}
\end{quote}

\textbf{Two pieces, left before right.} Count the pieces, not the strokes.

\textbf{The left piece, two strokes.} A vertical that runs down and flicks up to the right at its foot, then a short dot to the right of it, falling to the lower right.

\textbf{The right piece, two strokes.} A long left-falling stroke from the top right, sweeping down toward the foot of the first stroke, then a dot at the lower right, falling to the lower right.

2 + 2 = four.

\subsection*{Your turn}
\begin{itemize}
\raggedright
  \item \emph{Write it:} \zh{以}, all four, with the model in front of you
  \item \emph{Cover:} write it again from memory, and count the pieces, not the strokes
\end{itemize}

\begin{tcolorbox}[breakable,colback=teal!4,colframe=teal!35!black,title={Before you move on}]
How many strokes in \zh{以}? (\textbf{Four}.) Which stroke flicks up? (\textbf{The first, at its foot}.) Which way does the long one fall? (\textbf{Down to the left}.)

Source: \href{https://raw.githubusercontent.com/chanind/hanzi-writer-data/68d10a4b21150cae5e1ebbd223eed289cf32d90c/data/\%E4\%BB\%A5.json}{Hanzi Writer Data \zh{以}, pinned stroke paths and medians}.
\end{tcolorbox}
\section[gāngcái]{gāngcái — a moment ago}

\label{lesson:ZH-C126-gangcai}



\begin{quote}
Before the new one: say the Mandarin for sometimes, then the Mandarin for in the past.
\end{quote}

\subsection*{You'll want to know: gāngcái}
\emph{gāngcái} — \textquotedblleft{}a moment ago\textquotedblright{}.

\begin{quote}
\emph{gāngcái wǒ zài xǐshǒujiān} — A moment ago I was in the toilet.
\end{quote}

Its characters come later; for now it is a word for the ear.

\begin{cousinweb}
\emph{gāngcái} is built from \emph{gāng} and \emph{cái}, both inherited from Middle Chinese, \emph{kang} and \emph{dzoj} (Guangyun, in Baxter's transcription). The dictionary glosses the first 'strong', and lists the second, under an older way of writing it, as 'barely', a different word from the \emph{cái} of talent; how the pair came to mean a moment ago is not settled.
\end{cousinweb}

\begin{grammarlens}[title={What you've built}]
One more word for times of day.
\end{grammarlens}

\subsection*{Your turn}
\begin{itemize}
\raggedright
  \item \emph{Say it:} \emph{gāngcái}
  \item \emph{Say it:} \emph{gāngcái}, once more
  \item \emph{Say it:} \emph{gāngcái wǒ zài xǐshǒujiān} — A moment ago I was in the toilet.
  \item \emph{Recall:} say the Mandarin for sometimes, then the Mandarin for in the past, then say \emph{gāngcái} again
\end{itemize}

\begin{tcolorbox}[breakable,colback=teal!4,colframe=teal!35!black,title={Before you move on}]
What does gāngcái mean? (\textquotedblleft{}A moment ago\textquotedblright{}.) Say it once more.
\end{tcolorbox}
\section[yǐjīng]{yǐjīng — already}

\label{lesson:ZH-C126-yijing}



\begin{quote}
Before the new one: say the Mandarin for in the past, then the Mandarin for a moment ago.
\end{quote}

\subsection*{You'll want to know: yǐjīng}
\emph{yǐjīng} — \textquotedblleft{}already\textquotedblright{}.

\begin{quote}
\emph{wǒ yǐjīng chī fàn le} — I've already eaten.
\end{quote}

Its characters come later; for now it is a word for the ear.

\begin{cousinweb}
\emph{yǐjīng} is built from \emph{yǐ} and \emph{jīng}, both inherited from Middle Chinese, \emph{yiX} and \emph{keng} (Guangyun, in Baxter's transcription). The dictionary glosses the first 'to stop' and 'finished', a different word from the \zh{以} of \zh{以前}, and the second 'constant' and 'the warp of cloth'; how the pair came to mean already is not settled.
\end{cousinweb}

\begin{grammarlens}[title={What you've built}]
One more word for times of day.
\end{grammarlens}

\subsection*{Your turn}
\begin{itemize}
\raggedright
  \item \emph{Say it:} \emph{yǐjīng}
  \item \emph{Say it:} \emph{yǐjīng}, once more
  \item \emph{Say it:} \emph{wǒ yǐjīng chī fàn le} — I've already eaten.
  \item \emph{Recall:} say the Mandarin for in the past, then the Mandarin for a moment ago, then say \emph{yǐjīng} again
\end{itemize}

\begin{tcolorbox}[breakable,colback=teal!4,colframe=teal!35!black,title={Before you move on}]
What does yǐjīng mean? (\textquotedblleft{}Already\textquotedblright{}.) Say it once more.
\end{tcolorbox}
